\documentclass[11pt]{article}
\usepackage[a4paper,margin=25mm]{geometry}
\usepackage[T1]{fontenc}
\usepackage{lmodern,amsmath,hyperref,enumitem}
\hypersetup{hidelinks}
\setlength{\parindent}{0pt}
\setlength{\parskip}{6pt}
\newcommand{\comment}[1]{\textit{Comment: #1}}
\begin{document}
\begin{center}
{\Large Response to the editor and reviewers}\\[6pt]
Manuscript ARRAY-D-26-03529\\
Mathias Valla\\
Original title: Sparse farther-sighted trees can improve forest accuracy--time tradeoffs\\
Revised title: Bounded split optimization in decision trees and forests: gains, costs, and limits
\end{center}

Dear Dr Hazarika,

Thank you for the opportunity to revise the manuscript. I agree with Reviewer 1 that the original depth-three, single-split screening was insufficient to support its practical conclusions. The revision changes the evidence as well as the framing: it adds repeated stratified outer partitions, deeper and unrestricted tree capacities, training-only tuned conventional baselines, paired uncertainty estimates, and multiplicity-adjusted comparisons. It also corrects implementation and timing issues discovered during the revision. The contribution is now explicitly an empirical evaluation of an established induction rule, not a new lookahead algorithm.

Having considered your suggestion, I would now welcome the possibility of a transfer to another suitable journal offering publication without a mandatory article processing charge (APC), including a subscription publication route if available. I would be grateful for your guidance on an appropriate receiving journal and confirmation of its applicable publication terms before agreeing to a transfer. My preference is for a no-APC route rather than paid open access. I understand that a transfer remains subject to the receiving journal's independent assessment and does not imply acceptance.

Reviewer 2's report appears to concern a different manuscript: it describes a numerical-aware LogBERT extension for telecommunications monitoring, whereas this submission concerns decision-tree induction and bootstrap ensembles. I respectfully ask you to verify whether that report was attached to the correct submission. Below I identify the mismatch point by point, without attributing its cause or attempting to answer unrelated technical questions.

\section*{Response to the scientific handling editor}
\begin{enumerate}[label=E\arabic*.]
\item \comment{Significant revisions are required; another Elsevier journal may be considered.}

The methods and framing have been substantially rewritten. Sections 3--4 distinguish search horizon from tree capacity and describe the corrected protocol. The complete new fixed and matched-grid, training-only tuned evaluation replaces the preceding results in Section 5 and informs the discussion and conclusion in Sections 6--7. All five shards in both stages completed, and the full independent analysis passed. As stated above, I would welcome a suitable transfer with no mandatory APC and would appreciate confirmation of the publication terms before agreeing.

\item \comment{The author list must not change.}

The author list is unchanged. Mathias Valla remains the sole author and corresponding author; the affiliation and contact details are retained.

\item \comment{Suggested references should be assessed independently for relevance.}

The related-work section now discusses direct historical lookahead and forest precedents, contemporary scalable and hybrid tree search, and evidence that objectives and tuning affect comparisons. References are included because they support the decision-tree study. BARTPredict and MTRC are not included: the attached report does not establish their relevance to this manuscript.

\item \comment{Revision does not guarantee publication.}

I understand and appreciate the opportunity for reassessment. The revised manuscript makes no assertion that the changes establish publishability or guarantee acceptance.
\end{enumerate}

\section*{Response to Reviewer 1}
\subsection*{R1.1. Depth-three restriction and practical scalability}
\comment{The depth limit is unrealistic for practical random forests, and lookahead is computationally prohibitive at realistic depths.}

I agree that a depth-three comparison alone cannot establish practical competitiveness. The new protocol retains that controlled setting but additionally tests maximum final depth six and no depth limit on the same datasets and partitions. CART and one-, two-, and three-sighted trees have the same common hyperparameters and allowed total depth, with no post-pruning in either implementation. The fixed-architecture forest comparison also matches feature subsampling, bootstrap multiplicities, and ensemble size. A separate training-only selection experiment is applied to both conventional and mixed-capable forest families, rather than tuning only the baseline. Its shared structural grid includes deeper configurations; sharing that grid does not imply equal candidate counts or computation. Both stages are complete and their full analysis has passed validation.

The Introduction and the new related-work subsection ``Tree-size regularization'' clarify why a shallow regime is nevertheless scientifically relevant. Zhou and Mentch (2023), \emph{Trees, forests, chickens, and eggs: when and why to prune trees in a random forest}, provide direct evidence for tree-size regularization, particularly at low signal-to-noise ratios (\href{https://doi.org/10.1002/sam.11594}{doi:10.1002/sam.11594}). Duroux and Scornet (2018), \emph{Impact of subsampling and tree depth on random forests}, analyze early termination and subsampling in quantile regression forests and examine tree-size choices empirically (\href{https://doi.org/10.1051/ps/2018008}{doi:10.1051/ps/2018008}). These are conditional results, with regression, leaf-count, and resampling settings that differ from ours. They motivate studying shallow members but do not validate depth three as a general optimum for Gini classification forests. The cited consistency theory also concerns growing tree sizes, not a fixed eight-leaf limit.

Our depth-three design limits each member to at most eight leaves; it does not limit the aggregated forest to eight prediction regions or make a large ensemble as interpretable as one small tree. The revised manuscript acknowledges that deeper forests can reduce underfitting and improve prediction. Shallow-tree results from boosting are not used to justify bootstrap-forest performance. Thus this literature clarifies the research scope, while the added deep and tuned controls address the reviewer's empirical objection.

Section 3 now separates the local optimization horizon $k$ from the final tree-depth limit $D$. A two-generation search is restored at each actual node and can build a deeper tree; its horizon does not grow with $D$. I corrected the previous complexity discussion: complete split-assignment counts are not the work performed by the additive recursion. A uniform candidate-count abstraction gives $E_h\leq m+2mE_{h-1}$, before sorting and sample-partition costs. This still produces expensive search, but does not prove that every modern non-greedy method is unscalable.

The revised claims remain limited to the fixed PMLB sample. Unlimited final depth on small tasks is not evidence of scalability to large/high-dimensional modern workloads. Three-sighted members are now included in the completed repeated fixed and tuned experiments, rather than represented only by historical single-split screening. Relevant locations: Sections 3.1--3.3, 4.2, the depth results in Section 5 and Figure 1, and the limitations in Section 6.

\subsection*{R1.2. Single split, small differences, and lack of statistical analysis}
\comment{Without cross-validation, repeated runs, tuning, or significance tests, the observed gains cannot be separated from random variation.}

I agree. The new protocol uses five paired stratified outer partitions per dataset, with changing estimator/bootstrap seeds and an outer test fraction of one quarter. Fixed, unpruned CART and sighted models are compared with matching common settings. Separately, conventional RF, the one-/two-sighted family, and the inclusive one-/two-/three-sighted family all undergo inner cross-validation using the same outer-training rows and fold indices. All share the same depth, leaf-size, feature, and ensemble-size grid; both mixed families include the entire RF candidate space. Three configurations are locked before separate refits and outer-test evaluation. Thus tuned RF is compared with tuned mixtures, never presented as a fair contrast with an untuned mixture. Larger mixed candidate spaces are acknowledged and their full required selection work is reported, not hidden behind the cost of a cheap selected refit.

Inner cross-validation is stratified when every observed training class has at least two examples. A singleton or single-class training partition uses a documented ordinary shuffled three-fold fallback; this is not reliable rare-class validation. Ensemble sizes 20, 40, 60, and 100 are available to every tuned family. The 200-member fixed curves are descriptive and not tuning options available only to one competitor.

The balanced-accuracy sensitivity uses the library's mean recall over classes present in the scoring partition. A predicted class absent from that partition has undefined recall and is excluded from this mean; the library warns about this situation. This is not a fit failure, and the metric cannot establish recall for absent rare classes. The primary accuracy criterion and the common split protocol are unchanged.

The completed analysis first averages paired repeat effects within each dataset. It reports equal-dataset mean and median differences, wins/ties/losses, balanced-accuracy sensitivity, unadjusted percentile dataset-bootstrap intervals, and two-sided Wilcoxon tests with one Holm adjustment across eleven prespecified fixed and tuned contrasts. All other grid points are descriptive. Repeated overlapping holdouts are not counted as independent datasets. The bootstrap estimates a mean effect and the signed-rank test a location hypothesis, with its usual symmetry qualification; those are distinguished. Known related benchmark families and named constructed/synthetic tasks receive separate sensitivity analyses. A nonsignificant result is not interpreted as equivalence.

% BEGIN GENERATED response_results
% Replace the obsolete response_results status block; integrate the additional
% completion/tense corrections listed in full_interpretation_notes.md.
% First scientific writing pass, not a journal acceptance assessment.

The fixed and training-selected stages are now complete on all 57 datasets
and five paired outer partitions, with no missing or skipped task. The fixed
stage comprises 1,710 structural blocks, 5,130 single-tree scores, and 136,800
forest scores. The selected stage provides three separately refitted models
for each of 285 outer tasks, adding 855 scores. Independent full analysis
reconstructed stored fixed and inner-validation forest predictions and
constituent fitting costs and independently reconstructed the validation
winners. Single-tree and direct-refit scores received range and checkpoint
consistency checks, not independent prediction reconstruction. The complete
analysis now applies the single prespecified eleven-contrast Holm adjustment.
Figures 2--3 and Tables 1--2 use these completed records, rather than the
preceding tuned-versus-untuned evaluation.

The deeper controls materially qualify the shallow result. For all-feature
single trees, horizons two/three have descriptive mean differences from CART
of $2.77/3.54$ points at depth three, $0.39/0.59$ at depth six, and
$-5.13/-8.29$ without a depth limit. At 100 members, pure all-feature
horizon-two/three forests have mean differences of $1.90/2.34$ points
at depth three but $-6.84/-9.68$ without a depth limit. The square-root
feature setting differs: its unrestricted pure-forest mean differences
are $1.08/0.78$ points. We therefore do not generalize the all-feature
depth pattern to conventional feature-subsampled forests or claim that
shallow members are universally preferable. Selected RF chooses final
depth six or unrestricted depth in 219 of the 285 tasks. The conditional
tree-size literature motivates the shallow question, not depth three as
a general optimum, and boosting results are not treated as bootstrap-forest
evidence.

For the primary depth-three all-feature comparison, replacing ten of 100
CART members with horizon-two/three trees gives mean gains of $1.06$
[$0.45$, $1.77$] and $1.27$ [$0.61$, $2.02$] points. Both medians
are zero, with $27/21/9$ and $28/14/15$ wins/ties/losses, and joint
$p_H=0.00996/0.02066$. With square-root features, gains are $0.42/1.04$
points and $p_H=0.09539/0.00150$. The unrestricted replacement medians
are all zero and no corresponding adjusted signed-rank value is below
$0.05$. These facts are not interpreted as equivalence. Mean effects and
their pointwise bootstrap intervals are distinguished from the signed-rank
location hypothesis throughout.

Crucially, selected families are compared with selected RF. The inclusive
family's mean accuracy gain is $1.5693$ points [$0.2991$, $3.2767$],
with median zero, $28/11/18$ wins/ties/losses, and joint $p_H=0.48838$.
The one-/two-sighted family's gain is $0.9674$ points [$-0.0668$, $2.2979$],
with median $0.00617$ points, $29/8/20$ wins/ties/losses, and $p_H=1$.
The positive inclusive mean interval is not contradicted by its non-rejected
signed-rank location null; these summaries have different estimands. We do
not describe the result as no improvement or equivalence. However, the
inclusive mean gain becomes $1.03$ points [$-0.27$, $2.90$] under the
48-family weighting and $0.53$ [$-0.91$, $2.77$] on the 38 tasks outside
the named synthetic/constructed set. These intervals cross zero and limit a broad
generalization claim. Moreover, 133 of 285 inclusive selections are pure
forests and 152 are mixed (146 two-horizon and six three-horizon mixtures).
The selected-family effect cannot be attributed solely to sparse mixtures.

All families receive the same common structural and ensemble-size grid,
not equal computation: their candidate spaces contain 48, 288, and 768
configurations. Standalone cross-validation fitting work averages $2.271$,
$112.079$, and $2,671.049$ s for RF, one-/two-sighted, and inclusive
families, including banks not represented in the eventual selected model.
Direct refit wall times are $0.0626$, $0.9518$, and $12.9374$ s, including
construction/bootstrap preparation but excluding prediction. The sums of
measured workflow stage wall times are $3.405$, $118.035$, and $2,693.803$
s. These exclude checkpoint I/O, startup, downloads, and idle gaps and are
not uninterrupted end-to-end measurements. Physical shared-bank accounting
counts common work once, rather than summing the standalone family charges.
The much larger search costs remain a central limitation despite positive
estimated tuned mean effects.

Figure 2 is discussed immediately after the plot as a descriptive mean-cost
comparison. At the $0.7$ s mean-work reference, the highest observed mean
accuracy among plotted points below the line is the pure horizon-two,
20-member forest ($74.79\%$ at $0.3381$ s). Sparse mixtures do not
necessarily improve on smaller pure farther-sighted forests. The line is
not a per-dataset budget guarantee, and interpolation is not an evaluated
model. The ad hoc utility is not reinstated. The revised conclusion is a
conditional empirical assessment of established bounded search, not a new
algorithm, a causal diversity explanation, or a general accuracy--cost
superiority claim.
% END GENERATED response_results

Relevant locations: Sections 4.1--4.4 and 5, Tables 1--2, Figures 1--3, the complete CSV results, paired-comparison table, dataset manifest, and tuning records. The original single-split gains are not reused as uncertainty evidence.

\subsection*{R1.3. Limited novelty and no new search efficiency}
\comment{Lookahead induction is old, and the manuscript does not overcome its computational bottleneck.}

I agree and have changed the title, abstract, introduction, related work, and conclusion accordingly. Section 2 now treats Murthy and Salzberg's two-level LOOK rule as direct prior art, discusses stepwise lookahead forests, and positions the study against modern cached/dynamic-programming and hybrid methods, including DL8.5, MurTree, STreeD, and SPLIT. The recent literature on objective and tuning choices is also discussed.

The paper no longer presents bounded search or forest heterogeneity as a new algorithm. Its contribution is a reproducible assessment of when sparse horizon allocation changes prediction and fitting costs, and when its apparent benefits depend on the comparator or capacity. The isolated Cython scorer reuses terminal-split prefix class counts and improves numerical-search implementation efficiency; it does not change the underlying bounded optimization criterion or its worst-case horizon dependence. We do not claim a new non-greedy algorithm or asymptotic result. We have not performed a solver benchmark against those modern methods, and the text does not rank them by accuracy or efficiency.

\subsection*{R1.4. Ad hoc utility heuristic}
\comment{The accuracy--time utility is simplistic and ad hoc.}

I agree. I removed the logarithmic utility, its break-even coefficients, and the claim that a plotted vertical mean-time slice validates fixed-budget superiority. The revised study reports paired predictive differences and measured fitting costs directly. Figure 2 is explicitly descriptive: it does not select a model from outer-test scores or represent a validated per-dataset resource policy. Section 6 explains what a future training-only budget-selection experiment would require, including accounting for selection cost and checking final fitting feasibility.

\subsection*{Additional corrections identified during revision}
The revision also corrects three issues not explicitly raised in the report:
\begin{enumerate}
\item Historical mixtures joined prefixes of separately fitted forests and could reuse bootstrap positions. New mixtures replace distinct positions, with exactly one member at each position and correct class-probability alignment.
\item Historical mixed times were allocated from pure-bank averages. New controlled-ensemble costs sum the actual timings of the included members; selected-model whole-refit wall times are retained separately. Timing boundaries, excluded preparation/prediction, baseline selection cost, and concurrent workload are stated.
\item The preceding custom greedy learner differed from standard CART in zero-gain stopping and bootstrap sample-count semantics. The new protocol uses actual library CART as its one-sighted baseline, allows feasible zero-gain splits in sighted trees, and uses identical weighted bootstrap draws. A verification harness establishes RF-bank prediction agreement within absolute tolerance $10^{-12}$, with zero relative tolerance, at every considered ensemble size, and tests depth-matched interaction recovery without relying on a zero-gain stopping difference. Numerical and tie-breaking differences between implementations remain disclosed.
\end{enumerate}
These changes are reported in the manuscript and reproducibility documentation. Historical aggregates without retained repeat-level evidence are separated from the new complete evaluation and do not support its inference.

\section*{Response to Reviewer 2}
I appreciate the work involved in reviewing a manuscript. However, the supplied report's technical content does not correspond to ARRAY-D-26-03529. For clarity, its individual points are addressed as follows:
\begin{enumerate}[label=R2.\arabic*.]
\item \comment{A numerical-aware extension of LogBERT for telecom monitoring; positioning against BARTPredict and MTRC.}

This manuscript does not propose or evaluate LogBERT, a telecommunications system, or a proactive LLM framework. It studies Gini-based decision-tree split optimization and bootstrap forests. I request confirmation of the report's assignment.

\item \comment{Congestion ground truth, chronological splitting, normalization, threshold calibration, and numerical embedding are insufficiently specified.}

There is no congestion ground truth, numerical embedding, or anomaly threshold in the submitted study. Its labels are PMLB classification targets. The revised paper specifies stratified outer partitions and training-only inner model selection; a chronological telecom protocol would address a different research problem.

\item \comment{Check the anomaly-score formulation.}

No anomaly score is defined. The manuscript's criterion is weighted terminal Gini impurity, now specified by an additive recursion in Section 3.1.

\item \comment{Insufficient uncertainty, statistical analysis, per-class results, and evidence that the model predicts.}

Although the report appears to concern another model, uncertainty and statistical assessment are relevant general concerns and are addressed by the repeated, paired analysis described under R1.2. The revised study reports held-out classification accuracy and balanced accuracy. It makes no telecommunications forecasting or class-specific deployment claim; a telecom per-class validation cannot be supplied for this study.

\item \comment{Reduce or support claims about cyberattack detection, generalization, real-time performance, and production readiness.}

The manuscript makes no cyberattack, real-time monitoring, or production-readiness claim. Its generalization statements are confined to held-out performance on the stated benchmark sample. I would be grateful for any corrected report intended for this submission and would address it in a further response.
\end{enumerate}

Thank you for considering the revision and for clarifying the apparent report mismatch.

Yours sincerely,\\
Mathias Valla\\
Chaire DIALog, Institut Louis Bachelier, Paris, France\\
\href{mailto:mathias.valla@gmail.com}{mathias.valla@gmail.com}
\end{document}
